\chapter{Dodatni rezultati}
\label{app:dodatni-rezultati}

\TODO{
\begin{itemize}
    \item Sem dodaj dodatne tabele po foldih, rezultate manj pomembnih modelov in rezultate, ki niso šli v glavno besedilo.
    \item Dodaj rezultate za \textit{subject LOO} in \textit{recording LOO} pri vseh signalnih naborih: pogled, zenici, pogled+zenici in vsi signali.
    \item Dodaj vse izračunane metrike rezultatov, ne samo \textit{accuracy} in makro-F1, ki ostaneta v glavnem besedilu.
    \item Razširjene metrike: balanced accuracy, weighted-F1, AUC, po potrebi precision/recall po razredih in rezultati po posameznih gubah.
    \item Dodatne matrike zmede: recording LOO, binarne naloge, najboljši negrafovski model in primerjave, ki niso nujne za glavno zgodbo.
    \item Dodaj rezultate primerjave različnih tipov grafovskih konvolucij oziroma GNN slojev, če ostanejo zunaj glavnega besedila.
    \item Dodaj rezultate 3-razredne klasifikacije in kratek komentar primerjave z izvirnim člankom podatkovne zbirke MAHNOB-HCI oziroma HCI-tagging.
    \item Dodatne porazdelitve podatkov: porazdelitev razredov po subjektih, po posnetkih, po izvorni čustveni oznaki in za binarni low/high nalogi.
    \item Dodatne slike ujemanja oznak: podrobnejše heatmape Table 6 razredov proti samoocenam in povzetki label-noise analize, če v glavnem besedilu ostane samo ena skupna slika.
    \item UMAP procesiranih podatkov dodaj sem samo, če je slika zanimiva, vendar premalo ključna ali premalo čista za glavno besedilo. Če UMAP ne doda interpretativne vrednosti, ga ne vključuj.
\end{itemize}
}
